\chapter{Febrile Neutropenia}

\section*{Definition}
Febrile neutropenia is a fever that develops when you have an abnormally low number of neutrophils, a type of white blood cell that fights bacterial and fungal infection. It is most often a complication of cancer chemotherapy, but it can also follow stem-cell transplantation or occur with leukemia. It is a medical emergency: with so few neutrophils, an infection can spread into the bloodstream and become life-threatening within hours, often before producing any obvious signs of inflammation.

\section*{When to See a Doctor}
Contact your cancer team or go to the emergency department immediately if you are receiving chemotherapy and:
\begin{itemize}
  \item Your temperature reaches 38\textdegree{}C (100.4\textdegree{}F) or higher, even if you feel otherwise well
  \item You have chills, shaking (rigors), or a sudden feeling of being very unwell
  \item You have new cough, shortness of breath, abdominal pain, diarrhea, or painful urination
  \item You feel dizzy, faint, or confused, or your heart is racing (possible sepsis)
  \item You have an intravenous line or port and notice redness, pain, or tenderness at the site
\end{itemize}
Do not wait for other symptoms -- fever alone in a neutropenic patient requires urgent antibiotics.

\section*{How It Develops}
Neutrophils are your body's first responders against bacteria and fungi. They circulate in the blood, and when they detect an invading organism they travel to the site and engulf it. Chemotherapy kills dividing cells, including the bone-marrow cells that make neutrophils, so the neutrophil count falls to its lowest point about 7--14 days after a cycle of treatment. During this "nadir" period your defenses are stripped away, and even a mild mouth sore, a tiny skin break, or a colonized intravenous line can let bacteria through. Because there are so few neutrophils, the usual warning signs of infection -- pus, swelling, and tenderness -- may not appear. Fever is often the only signal that an infection is underway, and without rapid antibiotics it can progress to bloodstream infection (sepsis) and shock.

\section*{Causes}
\begin{itemize}
  \item Chemotherapy, which suppresses the bone marrow and lowers the neutrophil count
  \item Hematologic cancers such as leukemia, lymphoma, and myeloma, which crowd out normal marrow cells
  \item Stem-cell transplantation and the medicines used to suppress the immune system afterward
  \item Radiotherapy to large areas of bone marrow
  \item Some medicines that damage the marrow, such as certain immunosuppressants
  \item The underlying infection that the low counts allowed to take hold
\end{itemize}

\section*{Related Symptoms}
\begin{itemize}
  \item A single temperature of 38\textdegree{}C (100.4\textdegree{}F) or higher, or a temperature below 36\textdegree{}C (96.8\textdegree{}F)
  \item Chills and rigors (uncontrollable shaking)
  \item Fatigue, weakness, and muscle aches
  \item Mouth ulcers and sore throat
  \item Cough, shortness of breath, or chest pain
  \item Abdominal pain, nausea, vomiting, or diarrhea
  \item Redness, pain, or drainage around a catheter or port site
\end{itemize}

\section*{Medical Evaluation}
\begin{itemize}
  \item An immediate complete blood count to check the neutrophil count and confirm neutropenia.
  \item Blood cultures drawn from both a vein and your central line, plus urine and sputum cultures and a chest X-ray.
  \item A careful examination of the mouth, skin, sinuses, lungs, abdomen, and around any catheter or port for subtle signs of infection.
  \item Blood tests for kidney and liver function and inflammatory markers, and monitoring for the complications of sepsis.
  \item A review of your cancer treatment, its timing, and any antibiotic prophylaxis you take.
\end{itemize}

\section*{Treatment Options}
\begin{itemize}
  \item Empiric broad-spectrum intravenous antibiotics started immediately, before culture results come back.
  \item Continued antibiotics based on which organism is found and whether fever and counts improve; oral antibiotics may be an option for carefully selected low-risk patients.
  \item A growth factor (G-CSF) such as filgrastim to speed neutrophil recovery in selected patients.
  \item Antifungal therapy added if fever persists beyond several days despite antibiotics, since fungal infection is a growing risk as neutropenia continues.
  \item Intravenous fluids and blood-pressure support if sepsis develops, and oxygen for lung infection.
  \item Preventive planning for future cycles: G-CSF, antibiotic prophylaxis, and dose adjustments of chemotherapy to reduce the risk of another episode.
\end{itemize}


\section*{Self-Care and Home Management}
\begin{itemize}
  \item Do not try to manage this at home -- even with a single episode of fever, go straight to your cancer center or emergency department.
  \item Call ahead so your oncology team knows you are coming and can start antibiotics without delay.
  \item Drink fluids to avoid dehydration, and take your temperature correctly with a working thermometer if you have one.
  \item Avoid giving fever-reducing medicines before your evaluation if possible, since they can mask how high the fever goes.
  \item Follow your chemotherapy team's written instructions for what to do when a fever develops, and report the exact temperature and time.
\end{itemize}

\section*{References}
\begin{thebibliography}{99}
\bibitem{febrile_neutropenia:idsa} Freifeld AG, Bow EJ, Sepkowitz KA, et al. Clinical practice guideline for the use of antimicrobial agents in neutropenic patients with cancer: 2010 update by the Infectious Diseases Society of America. \textit{Clin Infect Dis}. 2011;52:e56--e93.
\bibitem{febrile_neutropenia:esmo} Klastersky J, de Naurois J, Rolston K, et al. Management of febrile neutropenia: ESMO Clinical Practice Guidelines. \textit{Ann Oncol}. 2016;27(suppl 5):v111--v118.
\bibitem{febrile_neutropenia:nccn} National Comprehensive Cancer Network. Prevention and treatment of cancer-related infections. NCCN Clinical Practice Guidelines in Oncology; accessed 2025.
\end{thebibliography}
